\documentclass[11pt]{article}

\usepackage[margin=1in]{geometry}
\usepackage{booktabs}
\usepackage{longtable}
\usepackage{array}
\usepackage{enumitem}
\usepackage{amsmath,amssymb}
\usepackage[colorlinks,linkcolor=black,citecolor=black,urlcolor=blue]{hyperref}
\usepackage[T1]{fontenc}

\setlist[itemize]{leftmargin=1.4em, itemsep=0.25em, topsep=0.3em, parsep=0pt}

\newcommand{\revpoint}[1]{\textit{#1}\par\nopagebreak\vspace{0.35em}}
\newcommand{\revhead}[1]{\subsection*{#1}}

\title{Response plan for ja-2026-12808m}
\author{Internal note (not the letter to Coley)}
\date{19 August 2026}

\begin{document}
\maketitle

\noindent
Coley rejected the manuscript; resubmission is invited by 13~February~2027.
Reviewer~1 asked for minor revision and put novelty in the top~5\%.
Reviewer~3 asked for major revision and put significance in the top~5\%.
Reviewer~2 said publish elsewhere, on impact, not on correctness.

We will not run every requested calculation.
From the co-author note we keep five jobs: a constructed out-of-distribution test,
one hydrogen-bonded system, one comparison with a modern ML potential,
one density above PBE, and a rewrite of the Abstract, the last paragraph of the
Introduction, and the first paragraph of the Discussion so the paper argues that
$n(\mathbf{r})$ is available at every MD step.

This note walks every numbered comment in reviewer order.
The italic sentence is their point.
The bullets are what we will do, or have already done.
The table at the end is the work list.

\section{Reviewer 1}

\revhead{R1.1 Hydrogen-bonded size extrapolation}
\revpoint{The systems are still simple; how does this work when hydrogen bonds matter?}
\begin{itemize}
  \item Train on water clusters $n=2$--$6$, test on $n=8,10,12$ (the protocol they named).
  \item Label the same geometries at PBE-D4 and $\omega$B97M-V so the H-bond figure
        and the hybrid table share a dataset. Labelling is running: PBE train is done,
        $\omega$B97M-V train is in flight.
  \item Skip NMA, methanol--water, and AcAc. Malonaldehyde covers the intramolecular
        proton-transfer motif they also listed.
\end{itemize}

\revhead{R1.2 Out-of-distribution conformations}
\revpoint{OOD conformations are missing, and a second MD trajectory is not proof of non-overlap.}
\begin{itemize}
  \item Built a ladder and measured it: in-distribution test (control), size hold-out
        (weak, and we say so), ice cutouts, low-density droplets, malonaldehyde
        proton-transfer scan.
  \item Separation: Bhattacharyya overlap plus nearest-neighbour distance in CV space,
        in units of the training set's own spacing. The figure is density error versus
        that distance.
  \item Ethanol high-$T$ / strain OOD stays as a cheap extra, not the panel we lead with.
\end{itemize}

\revhead{R1.3 Polythiophene length and Figure 6}
\revpoint{The polythiophene discussion is too long, and Figure 6 carries little.}
\begin{itemize}
  \item Shorten the section. Move most of Figure~6 to the SI. Keep one oligomer IR
        comparison in the main text. Manuscript only.
\end{itemize}

\revhead{R1.4 Figure 7}
\revpoint{Figure 7 is standard and belongs in the SI.}
\begin{itemize}
  \item Move it. Manuscript only.
\end{itemize}

\revhead{R1.5 Dataset documentation and train/test overlap}
\revpoint{SI S2.2 is thin, and ``a separate MD run'' does not prove the test set is non-overlapping.}
\begin{itemize}
  \item Print the RMSD numbers we have: resorcinol test frames sit inside the training
        basin (median min-RMSD 0.093~\AA, all 1000 below 0.2~\AA); ethanethiol and
        thiophene tests do not.
  \item If we find the force field, SHAKE, and trajectory length, we add them; if not,
        we say they were not recovered. The new water splits already have the same
        overlap treatment.
\end{itemize}

\revhead{R1.6 Duplicated timing}
\revpoint{The 755~ms versus 5~min timing appears twice.}
\begin{itemize}
  \item Keep one instance. Ethanol CPU: 1.62~s/step. GPU: 467~ms/step.
\end{itemize}

\revhead{R1.7 Unsupported dielectric claim}
\revpoint{The dielectric / hydrogen-bond-network sentence on p.~25 is not supported.}
\begin{itemize}
  \item Delete it. If the water-cluster results later support a narrower claim about
        H-bonds, that claim goes next to those results, not in the Discussion as a
        promise.
\end{itemize}

\section{Reviewer 2}

\revhead{R2.1 Hazra--Sanvito density-to-IR workflow}
\revpoint{Hazra, Patil, and Sanvito already go density to energy, forces, and IR.}
\begin{itemize}
  \item Cite \textit{J.\ Chem.\ Phys.}\ 2025, 163, 174104.
  \item Their density is a Jacobi--Legendre cluster expansion on uracil. Ours is an
        equivariant network with density fitting and a SAD correction, plus size, OOD,
        and H-bond tests they did not. Citation only.
\end{itemize}

\revhead{R2.2 Faster general IR methods}
\revpoint{MACE4IR, AIMNet-2, TranSpec, and AIQM are faster and work across chemical space.}
\begin{itemize}
  \item Compare transferable MLIPs and xTB on the molecules we already trained
        (expand Figure~2).
  \item We will not retrain on Crambin, paracetamol, or FGG. TranSpec maps SMILES to
        a spectrum; that is a different task and we say so.
\end{itemize}

\revhead{R2.3 Oligomer size extrapolation}
\revpoint{Training through the hexamer and testing longer oligomers is a weak extrapolation, and the error grows.}
\begin{itemize}
  \item Agree that this test is weak.
  \item Harder tests: ice, droplets, malonaldehyde scan. Size-OOD overlap is about
        0.20, above the 0.05 bar we set for ice and droplets, and we report that.
\end{itemize}

\revhead{R2.4 Companion network and energy conservation}
\revpoint{How does the companion network work, and is the AIMD energy conserving?}
\begin{itemize}
  \item Figure~1: $\{\mathbf{R}_I\}\to\mathrm{DenSNet}\to n(\mathbf{r})\to\mathrm{EnergyNet}\to E$,
        with $F_I=-\partial E/\partial\mathbf{R}_I$. Gradients pass through the density
        features.
  \item Quote NVE drift only after the energy head matches the 2024 checkpoint
        (\texttt{96w7KyGG}). Current 0.18--0.30~eV/ps numbers are from a mismatched
        head and stay out of the paper.
\end{itemize}

\revhead{R2.5 Who produces dipoles}
\revpoint{Page 4 implies the companion network produces electronic observables.}
\begin{itemize}
  \item Dipoles come from the density (\texttt{int1e\_r}). The companion head outputs
        energy and forces. Fix page~4.
\end{itemize}

\revhead{R2.6 One framework or two networks}
\revpoint{The text says one framework; there are two networks.}
\begin{itemize}
  \item Both modules run at every MD step. IR does not need a post-processed trajectory
        if the dipole integrator is on. Drop ``single network''.
\end{itemize}

\revhead{R2.7 Electron-count normalization versus delta-learning}
\revpoint{If normalizing by electron count already weighted core and valence, why delta-learning?}
\begin{itemize}
  \item Normalization is a loss weight so raw grid errors in the core do not dominate.
        SAD changes the target to $\rho-\rho_{\mathrm{SAD}}$. They are not substitutes.
        Keep both.
\end{itemize}

\revhead{R2.8 Softplus and signed $\Delta n$}
\revpoint{Delta-density is negative in places; how does that sit with a softplus?}
\begin{itemize}
  \item Softplus is applied so the reconstructed total density stays non-negative.
        $\Delta n$ itself is signed.
  \item About half of the DF coefficients on water and ethanethiol tests are already
        negative (46.7\% and 51.0\%).
\end{itemize}

\revhead{R2.9 AIMD length}
\revpoint{State the AIMD length in the main text.}
\begin{itemize}
  \item Ethanol AIMD on disk is 10~ps. Put every trajectory length in the main text
        (same job as R3.6).
\end{itemize}

\revhead{R2.10 Density-error units}
\revpoint{Density errors have no units.}
\begin{itemize}
  \item MAE in $e/a_0^3$ and relative MAE $\int|\Delta\rho|\,\mathrm{d}V/N_e$.
  \item We will not invent a density analogue of 1~kcal/mol. Tie density error to
        dipole error and then to IR intensity.
\end{itemize}

\revhead{R2.11 Equilibrium structures}
\revpoint{The manuscript claims DFT-like equilibria without numbers.}
\begin{itemize}
  \item After the energy head is restored, report geometry RMSD versus DFT on a few
        small molecules.
  \item Do not quote the current ethanol energy of ${\sim}0.24$~eV; that is an
        architecture mismatch.
\end{itemize}

\revhead{R2.12 Missing density ML citation}
\revpoint{Cuevas-Zuvir\'ia and Pacios, \textit{J.\ Chem.\ Inf.\ Model.}\ 2020, is missing.}
\begin{itemize}
  \item Add it to the introduction.
\end{itemize}

\revhead{R2.13 Corresponding authors}
\revpoint{Corresponding authors disagree between the manuscript and the SI.}
\begin{itemize}
  \item Make them match. Confirm the list with Tuckerman, Burke, and M\"uller before
        resubmission.
\end{itemize}

\section{Reviewer 3}

\revhead{R3.1 Hybrid DFT densities}
\revpoint{All reference densities are PBE; try PBE0+MBD or $\omega$B97M-V/def2-TZVPD.}
\begin{itemize}
  \item Nothing in the architecture is PBE-specific.
  \item Label the same water and malonaldehyde geometries at $\omega$B97M-V/def2-TZVPD
        and at PBE-D4; one table. No hybrid MD, no hybrid IR, no retraining of ethanol
        or thiophene at hybrid.
\end{itemize}

\revhead{R3.2 Comparisons with ML force fields}
\revpoint{Compare SO3LR, MACE-OFF, GFN-xTB, DFTB, and the Crambin / paracetamol / FGG IR suites.}
\begin{itemize}
  \item Expand Figure~2 on our molecules.
  \item Decline Crambin, paracetamol, and FGG: those papers test transferable MLIPs,
        and DenSNet is fitted per system.
\end{itemize}

\revhead{R3.3 Interpreting density errors}
\revpoint{A density MAE of order $10^{-3}$ is hard to interpret; add relative errors and uncertainty.}
\begin{itemize}
  \item Relative MAE, plus density error against dipole error and IR intensity error.
  \item A three-seed ensemble on water if it is cheap. Uncertainty quantification is
        not a separate project.
\end{itemize}

\revhead{R3.4 Cost and magnitude of delta-learning}
\revpoint{Delta-learning needs a cost and a magnitude, not just a lower MAE.}
\begin{itemize}
  \item Time the SAD lookup. Histogram $|\Delta\rho|/\rho$.
  \item Train direct and delta models on water and compare them with relative errors
        (configs already exist).
\end{itemize}

\revhead{R3.5 Hyperparameters and cutoff}
\revpoint{Architecture parameters look un-optimized, and cutoff is known to matter.}
\begin{itemize}
  \item Write down the paper cutoffs (5.0~\AA\ thiophene, 15~Bohr small molecules).
  \item Sweep cutoff $\{4,5,6,8\}$~\AA\ on water at PBE only. No grid search.
\end{itemize}

\revhead{R3.6 IR convergence with trajectory length}
\revpoint{``Sufficient duration'' is not a number.}
\begin{itemize}
  \item State the lengths. Plot IR versus 50, 100, 200, 500~ps on ethanol and
        thiophene 2-mer (those jobs exist). Report peak position and integrated area
        as a function of length.
\end{itemize}

\revhead{R3.7 General-purpose density model}
\revpoint{Could DenSNet become a general-purpose density model, perhaps with active learning?}
\begin{itemize}
  \item Discussion only. Present models are system-specific. A later path is large
        density datasets, pretraining, fine-tuning, and active learning.
  \item We will not build that model in this revision. OMol/CSH can be mentioned as
        work already running, not as a result of the paper.
\end{itemize}

\revhead{R3.8 Spectral language}
\revpoint{Peak positions agree better than intensities; stop saying the spectra match across the full range.}
\begin{itemize}
  \item Peak positions in named windows. Intensities and areas agree less evenly.
        The recovered \TeX\ already has this wording; keep it.
\end{itemize}

\revhead{R3.m1 Two-stage training cost}
\revpoint{Training cost for the two-stage procedure, especially on the larger systems.}
\begin{itemize}
  \item GPU-hours in the SI.
\end{itemize}

\revhead{R3.m2 Architecture in Figure 1}
\revpoint{Figure 1 should include the architecture, not only the workflow.}
\begin{itemize}
  \item Same figure as R2.4.
\end{itemize}

\revhead{R3.m3 Other density-derived properties}
\revpoint{Which other electronic properties could come from the predicted density?}
\begin{itemize}
  \item Hirshfeld charges, ESP, polarizability: list as future observables. No new
        calculations.
\end{itemize}

\section{To-do}

\noindent
Columns: ID, task, kind, status, depends on.

{\small
\setlength{\tabcolsep}{4pt}
\begin{longtable}{
  >{\raggedright\arraybackslash}p{0.10\textwidth}
  >{\raggedright\arraybackslash}p{0.36\textwidth}
  >{\raggedright\arraybackslash}p{0.11\textwidth}
  >{\raggedright\arraybackslash}p{0.15\textwidth}
  >{\raggedright\arraybackslash}p{0.16\textwidth}}
\caption{Work list for the resubmission.}\\
\toprule
ID & Task & Kind & Status & Depends on \\
\midrule
\endfirsthead
\midrule
ID & Task & Kind & Status & Depends on \\
\midrule
\endhead
\bottomrule
\endlastfoot

\multicolumn{5}{l}{\emph{Compute}} \\
\midrule
R1.1 & Water PBE labels, train, ID and size-OOD metrics & DFT & running & -- \\
R3.1 & Water $\omega$B97M-V labels, same-geometry hybrid table & DFT & running & R1.1 geometries \\
R1.2, R2.3 & Ice, droplet, malonaldehyde labels, then error versus distance & DFT & running & train checkpoint \\
(ops) & Watchdog running so training actually starts & ops & down & -- \\
R3.5 & Cutoff sweep at PBE & analysis & waiting & water PBE train \\
R2.4, R2.11 & Restore energy head, then NVE and geo-opt RMSD & MD & blocked & 2024 checkpoint \\
R3.6 & Ethanol and thiophene IR versus length & MD & jobs exist & trajectories \\
R2.2, R3.2 & Figure~2 GPU timings versus MACE-OFF, AIMNet2, xTB & MD & partial & GPU QOS \\
R2.10, R3.3, R3.4 & Relative density error, $\Delta\rho$ histograms, SAD timing & analysis & scripts ready & trained water model \\
\midrule
\multicolumn{5}{l}{\emph{Manuscript (can start now)}} \\
\midrule
-- & Abstract, last Intro paragraph, first Discussion paragraph: $n(\mathbf{r})$ at every MD step & manuscript & not started & -- \\
R1.3, R1.4 & Shorten polythiophene; move Figures~6 and~7 & manuscript & outlined & -- \\
R1.7, R3.8 & Delete dielectric sentence; keep softer spectral wording & manuscript & TeX already edited & -- \\
R2.4--R2.6, R3.m2 & Architecture in Figure~1; dipoles versus energy head; two modules & manuscript & outlined & -- \\
R1.5, R2.9, R2.10 & SI S2.2, overlap table, AIMD lengths, units & manuscript & overlap computed & -- \\
R2.1, R2.12, R2.2 & Cite Hazra--Sanvito, Cuevas-Zuvir\'ia, Pracht, TranSpec, AIQM & manuscript & not started & -- \\
R2.13, R3.m1 & Corresponding authors; two-stage GPU-hours in SI & manuscript & not started & names from PIs \\
\midrule
\multicolumn{5}{l}{\emph{Decline}} \\
\midrule
R3.2 & Crambin, paracetamol, FGG & decline & -- & -- \\
R3.1 & Retrain ethanol, thiophene, or polythiophene at hybrid & decline & -- & -- \\
R3.7 & Universal DenSNet or active-learning campaign as a paper result & decline & -- & -- \\
R3.5 & Full hyperparameter search & decline & -- & -- \\
R1.3 & New polythiophene experiment & decline & -- & -- \\
\end{longtable}
}

\end{document}
